 \part{ブール代数}
	\section{双対原理}
		\begin{proof}
			\quad\par
			論理式$P(\bm{x})$の双対$P(\bm{x})^\mathrm{d}$は$\overline{P(\overline{\bm{x}})}$で表せる。
			$P(\bm{x})=Q(\bm{x})$ならば$\overline{P(\overline{\bm{x}})} = \overline{Q(\overline{\bm{x}})}$なので$P(\bm{x})^\mathrm{d} = Q(\bm{x})^\mathrm{d}$である。
		\end{proof}
	\section{応用}
		\subsection{2進Grayコードへの変換}
			$x$の2進表現からGrayコードへの変換は$x$の2進表現と、それを1ビット右シフトして先頭に0をつけたものとの排他的論理和で得られる。
			これで上手くいく理由を考えてみる。
			$x$の$n$bit 2進表現を$b_{n-1},b_{n-2},\cdots,b_1,b_0$とする。
			Grayコードへの変換操作を乱暴に言うと、$b_{n-1}$から$b_0$に向かって見ながら、$b_{n-1}$はそのままで、以降は1と0が交互に現れる箇所を1、そうでない箇所を0に置き換える。
			この操作で$x$と$x+1$のGrayコードのハミング距離が1になることを示す。
			\\\newline
			(1) 最下位桁が0である場合
			\par
			\[
				\begin{matrix}
					x &= b_{n-1},b_{n-2},\cdots,b_1,0   & \to & b_{n-1}, b_{n-1}\oplus b_{n-2},\cdots,b_2\oplus b_1, \textcolor{red}{b_1} \\
					& &\text{変換} &\\
					x+1 &= b_{n-1},b_{n-2},\cdots,b_1,1 & \to & b_{n-1}, b_{n-1}\oplus b_{n-2},\cdots,b_2\oplus b_1, \textcolor{red}{\overline{b_1}}
				\end{matrix}
			\]
			両者のGrayコードは最下位桁のみが異なっている。
			\\\newline
			(2) 最下位桁が1である場合
			\par
			$x$の下位$k$桁が1の連続であるとする。
			\[
				\begin{matrix}
					x & = b_{n-1},b_{n-2},\cdots,b_{k+1},0,\underbrace{1,\cdots,1}_{k\text{個}}   & \to & b_{n-1}, b_{n-1}\oplus b_{n-2},\cdots,b_{k+2}\oplus b_{k+1},\textcolor{red}{\overline{b_{k+1}}},1,\underbrace{0,\cdots,0}_{k-1\text{個}} \\
					& &\text{変換} &\\
					x+1 & = b_{n-1},b_{n-2},\cdots,b_{k+1},1,\underbrace{0,\cdots,0}_{k\text{個}} & \to & b_{n-1}, b_{n-1}\oplus b_{n-2},\cdots,b_{k+2}\oplus b_{k+1},\textcolor{red}{b_{k+1}},1,\underbrace{0,\cdots,0}_{k-1\text{個}}
				\end{matrix}
			\]
			両者のGrayコードは第$k+1$桁のみが異なっている。